\documentclass[11pt,a4paper]{article}
\usepackage[margin=1in]{geometry}
\usepackage{booktabs}
\usepackage{hyperref}
\usepackage{enumitem}

\title{A Practical 20-Minute Security Baseline Chain for Node.js Services}
\author{Security Claw Project}
\date{\today}

\begin{document}
\maketitle

\begin{abstract}
This document summarizes a practical security analysis workflow designed for Node.js services such as OpenClaw and NanoClaw. The workflow prioritizes common, widely adopted tools and enforces a strict 20-minute execution budget. We describe the chain design, implementation hardening decisions, and measured outcomes from local experiments.
\end{abstract}

\section{Problem Statement}
Many teams need a reproducible baseline security check that is fast enough to run in pull requests and frequent CI jobs. Our objective is to build a minimal but effective scan chain that can:
\begin{itemize}[leftmargin=1.5em]
  \item detect dependency vulnerabilities,
  \item detect risky code patterns,
  \item detect container and repository misconfigurations,
  \item complete within 1200 seconds.
\end{itemize}

\section{Security Chain Design}
The baseline chain consists of three sequential stages:
\begin{enumerate}[leftmargin=1.5em]
  \item \textbf{Dependency scan}: \texttt{npm audit --audit-level=high --json}
  \item \textbf{SAST scan}: \texttt{semgrep scan --config p/default --json}
  \item \textbf{Container and misconfiguration scan}: \texttt{trivy fs --scanners vuln,misconfig --severity HIGH,CRITICAL --format json}
\end{enumerate}

To support reuse across similar services, stage commands are configured through service profiles (JSON) and executed by a shared engine script.

\section{Implementation and Hardening}
The framework includes:
\begin{itemize}[leftmargin=1.5em]
  \item service profiles under \texttt{services/*.json},
  \item execution scripts under \texttt{engine/},
  \item CI workflows under \texttt{.github/workflows/}.
\end{itemize}

Key hardening decisions:
\begin{itemize}[leftmargin=1.5em]
  \item replace string-based shell execution with structured \texttt{command + args} invocation,
  \item avoid direct unsafe interpolation in workflow \texttt{run} blocks,
  \item preserve before/after comparison for patch validation through summary diffs.
\end{itemize}

\section{Experimental Setup}
\textbf{Target}: NanoClaw local repository.\\
\textbf{Budget}: 1200 seconds.\\
\textbf{Outputs}: JSON reports and summarized metrics.

\section{Measured Results}
Two representative runs were recorded:

\begin{table}[h]
\centering
\begin{tabular}{lcccc}
\toprule
Run Label & Elapsed (s) & npm audit (High/Critical) & Semgrep Findings & Trivy High/Critical \\
\midrule
exp2 & 7 & 0 / 0 & 49 & 3 \\
local\_framework\_test & 21 & 0 / 0 & 54 & 3 (+1 misconfig) \\
\bottomrule
\end{tabular}
\caption{Observed baseline chain outcomes under the 20-minute budget.}
\end{table}

All observed runs completed far below the 1200-second limit.

\section{Patch Validation Flow}
For vulnerability remediation, the process includes:
\begin{enumerate}[leftmargin=1.5em]
  \item baseline scan and summary snapshot,
  \item patch application,
  \item re-scan with the same chain,
  \item summary diff and policy gate check.
\end{enumerate}

A patch is considered successful when targeted findings disappear in the after-state and no new High/Critical issues are introduced.

\section{Limitations}
\begin{itemize}[leftmargin=1.5em]
  \item Semgrep findings include both true positives and rule-noise; triage is required.
  \item Vulnerability presence in dependencies may require upstream package updates, not only local code changes.
  \item Results can vary with tool database update time.
\end{itemize}

\section{Conclusion}
This work demonstrates that a minimal, reusable, and hardened security chain can provide meaningful baseline coverage for Node.js services within strict CI time constraints. The approach is practical for continuous monitoring and can be extended with service-specific policy gates and patch regression checks.

\end{document}
